\chapter{Prosjektgjennomføring}
\label{ch:process}
Utviklingen har gått fra et sekvensielt prototypeprogram og usikre
sensorobservasjoner til et mer avgrenset system med testbar kontrollflyt,
tydelig feilrespons og sporbare datakilder. Daterte logger dokumenterer
oppstartsregresjon og sensorarbeid 21.--23. oktober, 60 cm-referansen
23. oktober samt Pi-integrasjon, exceptions, lys og adaptivt søk 2. november.
Den tidligere lange rapporten er bevart som \kode{development-history.tex};
daterte resultater er ikke omskrevet til nye resultater.

Gjeldende arbeid ligger i
\url{https://github.com/kidi970/ADAptive-VE3140}, på arbeidsgrenen
\kode{codex/pi-integration-exceptions-2026-11-02}.
Undermodulen Ada Drivers Library er låst til
\kode{8e90f85575666684cc6e0b51830494de259bd3bc}; den kumulative
\kode{patches/ada-drivers-2026-10-23.patch} må anvendes én gang etter
uthenting. Lokalt er denne patchen allerede anvendt. Verktøyversjoner,
prosjektfiler og hashmanifest følger dokumentasjonen.

For hver programvareendring brukes en avgrenset begrunnelse og relevante
kontroller. Feiltester skiller mellom applikasjonen og eksplisitte fakes;
mutasjoner demonstrerer at nye testporter faktisk kan avvise feil.
GitHub-oppsettet bygger prosjektet og tester kode ved relevante endringer;
lokale resultater dokumenteres uavhengig av om en ekstern CI-kjøring er utført.
Rapportens kodevedlegg kopieres automatisk fra gjeldende kilde og
kontrolleres byte for byte for å unngå utdaterte kodeutdrag.

Brukeren oppga omtrent 2--3 uker til demonstrasjon og én måned til eksamen.
Eksakte datoer er ikke bekreftet. Videre arbeid prioriteres etter avhengighet:
lys og strøm, sensorer og I2C, løftede hjul, tidsmåling og så løype.
Fysisk sensor-/bussfunksjon er den viktigste åpne avhengigheten.
Mer avansert regulering utsettes til nødvendig tilbakekobling og måledata finnes.

\begin{tabularx}{\textwidth}{@{}P{42mm}Y@{}}\toprule
Gruppemedlem & Må fylles av gruppen\\\midrule
Nicolay Olafsen, 266584 & \todo{Eget ansvar, konkrete bidrag, gjennomganger og læringsutbytte.}\\
Jacob Hellumstand Utengen, 228979 & \todo{Eget ansvar, konkrete bidrag, gjennomganger og læringsutbytte.}\\\bottomrule
\end{tabularx}

\todo{Dokumenter faktisk samarbeid, plan mot gjennomføring og hvilke valg
gruppen selv tok. Avstem den endelige KI-deklarasjonen og presentasjonskravene.}
Disse opplysningene kan ikke ferdigstilles fra kode eller oppdiktes i rapporten.
